\textbf{Model.} We use \texttt{facebook/convnext-tiny-224}
(ConvNeXt-Tiny), pretrained on ImageNet-1k and served through the Hugging
Face \texttt{transformers} library. All images are resized to
$224\times224$ and passed through the model's standard image processor
before inference.

\textbf{SHAP explainer.} Because ConvNeXt-Tiny is not a simple
gradient-friendly network in the sense required by SHAP's gradient-based
explainers, we use \texttt{shap.Explainer} with a model-agnostic partition
explainer over an \texttt{Image} masker configured with Gaussian blur
(\texttt{blur(32,32)}). The masker perturbs the image by blurring out
candidate superpixel regions and observes the resulting change in the
model's softmax output for the predicted class; SHAP then attributes
credit to each region based on its marginal contribution across many such
maskings. We cap the explainer at \texttt{max\_evals=100} per image, which
keeps runtime to roughly 30 seconds per image on CPU while still producing
a stable coarse-grained attribution map.

\textbf{Test set.} We selected 10 groups of ImageNet classes that are
visually similar enough to be plausibly confused by a classifier and by
a human at a glance: two dog breeds, three cat variants, two bird species,
two shark species, two big cats, three fox species, two elephant species,
two bear species, two terrier breeds, and two retriever breeds (22 images
total, one representative image per class, drawn from a public ImageNet
sample set). Image acquisition is scripted in
\texttt{code/fetch\_images.py} so the set can be extended by adding entries
to the class-pair table without changing the analysis code.

\textbf{Pipeline.} For each image we (1) run the classifier and record the
top-1 predicted label and confidence, (2) compute a SHAP attribution map
for that predicted label, and (3) save the resulting heatmap alongside the
source image. This is implemented in \texttt{code/run\_shap.py} and driven
end-to-end by \texttt{make loop} (fetch images, run SHAP, compile the
paper).
